\section{Evaluation Metrics}
\label{sec:performance_metrics}
%%%%%%%%%%%%%%%%%%%%%%%%%%%%%%%%%%%%%%%%%%%%%%%%%%%%%%%%%%%%%%%%%%%%%
\begin{table*}[t]
\centering
\resizebox{0.70\textwidth}{!}{
\begin{tabular}{l lcccc}
\toprule
\textbf{Pipeline} & \textbf{Partition} & \textbf{Accuracy} & \textbf{Precision} & \textbf{Recall} & \textbf{F1-score} \\
\midrule
\multirow{2}{*}
{CaseText Only}
  & Single & 62.27 & 33.50 & 30.88 & 29.45 \\
  & Multi  & 53.10 & 25.26 & 23.95 & 20.81 \\
\midrule
\multirow{2}{*}
{CaseText + Statutes}
  & Single & \textbf{67.07} & 45.29 & 44.55 & 44.32 \\
  & Multi  & 60.36 & \textbf{64.22} & \textbf{64.04} & 60.35 \\
\midrule
\multirow{2}{*}
{CaseText + Precedents}
  & Single & 61.73 & 41.92 & 41.35 & 40.81 \\
  & Multi  & 57.53 & 61.34 & 61.19 & 57.53 \\
\midrule
\multirow{2}{*}
{CaseText + Previous Similar Cases}
  & Single & 57.53 & \textbf{61.34} & \textbf{61.19} & \textbf{57.53} \\
  & Multi  & 61.73 & 41.92 & 41.35 & 57.53 \\
\midrule
\multirow{2}{*}
{CaseText + Statutes + Precedents}
  & Single & 64.71 & 43.50 & 42.98 & 42.78 \\
  & Multi  & \textbf{65.86} & 63.94 & 63.99 & \textbf{63.96} \\
\midrule
\multirow{2}{*}
{CaseFacts Only}
  & Single & 51.13 & 51.36 & 51.30 & 50.68 \\
  & Multi  & 53.71 & 51.18 & 51.18 & 51.18 \\
\midrule
\multirow{2}{*}
{Facts + Statutes + Precedents}
  & Single & 50.58 & 33.57 & 33.56 & 33.24 \\
  & Multi  & 52.57 & 52.01 & 52.01 & 52.01 \\
\bottomrule
\end{tabular}
}
\vspace*{-2mm}
\caption{Performance of various pipelines on Binary and Multi-label LJP task (best results in bold)}
\label{tab:prediction_pipeline_results}
\end{table*}

%%%%%%%%%%%%%%%%%%%%%%%%%%%%%%%%%%%%%%%%%%%%%%%%%%%%%%%%%%%%%%%%%%%%%
To evaluate the effectiveness of our Retrieval-Augmented Legal Judgment Prediction framework, we adopt a comprehensive set of metrics covering both classification accuracy and explanation quality. The evaluation is conducted on two fronts: the judgment prediction task and the explanation generation task. These metrics are selected to ensure a holistic assessment of model performance in the legal domain. We report Macro Precision, Macro Recall, Macro F1, and Accuracy for judgment prediction, and we use both quantitative and qualitative methods to evaluate the quality of explanations generated by the model.

\begin{enumerate}
    \item \textbf{Lexical-based Evaluation:} We utilized standard lexical similarity metrics, including Rouge-L~\cite{lin-2004-rouge}, BLEU \cite{papineni-etal-2002-bleu}, and METEOR \cite{banerjee-lavie-2005-meteor}. These metrics measure the overlap and order of words between the generated explanations and the reference texts, providing a quantitative assessment of the lexical accuracy of the model outputs.

    \item \textbf{Semantic Similarity-based Evaluation:} To capture the semantic quality of the generated explanations, we employed BERTScore \cite{BERTScore}, which measures the semantic similarity between the generated text and the reference explanations. Additionally, we used BLANC \cite{blanc}, a metric that estimates the quality of generated text without a gold standard, to evaluate the model's ability to produce semantically meaningful and contextually relevant explanations.

    \item \textbf{LLM-based Evaluation (LLM-as-a-Judge):} To complement traditional metrics, we incorporate an automatic evaluation strategy that uses large language models themselves as evaluators, commonly referred to as \textit{LLM-as-a-Judge}. This evaluation is crucial for assessing structured argumentation and legal correctness in a format aligned with expert judicial reasoning. We adopt G-Eval~\cite{liu-etal-2023-g}, a GPT-4-based evaluation framework tailored for natural language generation tasks. G-Eval leverages chain-of-thought prompting and structured scoring to assess explanations along three key criteria: \textit{factual accuracy}, \textit{completeness \& coverage}, and \textit{clarity \& coherence}. Each generated legal explanation is scored on a scale from 1 to 10 based on how well it aligns with the expected content and a reference document. The exact prompt format used for evaluation is shown in Appendix Table~\ref{tab:G-Eval Prompt}. For our experiments, we use the GPT-4o-mini model to generate reliable scores without manual intervention. This setup provides an interpretable, unified judgment metric that captures legal soundness, completeness of reasoning, and logical coherence, beyond what traditional similarity-based metrics can offer.

    % \item \textbf{Expert Evaluation:} To validate the interpretability and legal soundness of the generated explanations, we conducted an expert evaluation with three experienced legal professionals. A representative subset of outputs was rated on a 1–10 Likert scale, considering three criteria: factual accuracy, legal relevance, and completeness of reasoning. A score of 1 denoted poor or misleading explanations, while a score of 10 indicated high fidelity to legal standards and argumentative soundness. Each expert rated independently and was blind to the system configurations, thereby minimizing bias and ensuring objective assessment. This evaluation offered important insights beyond automated metrics, allowing us to better assess the practical utility and reliability of the proposed framework.


    \item \textbf{Expert Evaluation:} To validate the interpretability and legal soundness of the model-generated explanations, we conduct an expert evaluation involving legal professionals. They rate a representative subset of the generated outputs on a 1–10 Likert scale across three criteria: factual accuracy, legal relevance, and completeness of reasoning. A score of 1 denotes a poor or misleading explanation, while a 10 reflects high legal fidelity and argumentative soundness. This evaluation provides critical insights beyond automated metrics.

    \item \textbf{Inter-Annotator Agreement (IAA):} To further ensure the reliability and consistency of expert judgments, we computed multiple standard inter-rater agreement metrics. Specifically, Fleiss’ Kappa~\cite{fleiss1971measuring} was used to evaluate the overall agreement among multiple raters, Cohen’s Kappa~\cite{cohen1960coefficient} captured pairwise agreement while adjusting for chance, and the Intraclass Correlation Coefficient (ICC)~\cite{shrout1979intraclass} measured the reliability of continuous ratings across raters. In addition, Krippendorff’s Alpha~\cite{krippendorff2018content} provided a robust measure suitable for ordinal scales and missing data, while the Pearson Correlation Coefficient~\cite{benesty2009pearson} quantified the linear consistency of expert scores. The results demonstrated substantial agreement across different measures, reinforcing the credibility of the evaluation and lending strong support to the robustness of our findings.
\end{enumerate}

Together, these metrics provide a comprehensive, multi-perspective evaluation of our system’s ability to both predict judicial outcomes and generate legally coherent, interpretable rationales grounded in retrieved context.


%%%%%%%%%%%%%%%%%%%%%%%%%%%%%%%%%%%%%%%%%%%%%%%%%%%%%%%%%%%%%%%%%%%%%%%%%%%%%%%%%%%%%%%

\begin{table*}[t]
\centering
\resizebox{0.80\linewidth}{!}{
\begin{tabular}{lccccccc}
\toprule
\textbf{Pipelines} & \textbf{RL} & \textbf{BLEU} & \textbf{METEOR} & \textbf{BERTScore} & \textbf{BLANC} & \textbf{G-Eval} & \textbf{Expert Score} \\
\midrule
\multicolumn{8}{c}{\textbf{Single Partition}} \\
\midrule
CaseText Only & 0.16 & 0.03 & 0.18 & 0.52 & 0.08 & 4.17 & 5.2 \\
CaseText + Statutes & \textbf{0.17} & \textbf{0.03} & \textbf{0.20} & \textbf{0.53} & \textbf{0.09} & \textbf{4.21} & \textbf{5.5} \\
CaseText + Precedents & 0.16 & 0.03 & 0.19 & 0.51 & 0.08 & 3.45 & 4.6 \\
CaseText + Previous Similar Cases & 0.16 & 0.03 & 0.20 & 0.52 & 0.08 & 3.72 & 4.9 \\
CaseText + Statutes + Precedents & 0.16 & 0.03 & 0.19 & 0.52 & 0.08 & 4.11 & 5.4 \\
CaseFacts Only & 0.16 & 0.02 & 0.18 & 0.52 & 0.06 & 3.53 & 4.5 \\
Facts + Statutes + Precedents & 0.16 & 0.02 & 0.18 & 0.51 & 0.06 & 2.97 & 3.9 \\
\midrule
\multicolumn{8}{c}{\textbf{Multi Partition}} \\
\midrule
CaseText Only & 0.16 & 0.03 & 0.18 & 0.52 & 0.08 & 4.00 & 5.0 \\
CaseText + Statutes & \textbf{0.17} & \textbf{0.03} & \textbf{0.20} & \textbf{0.53} & \textbf{0.09} & \textbf{4.10} & \textbf{5.3} \\
CaseText + Precedents & 0.16 & 0.03 & 0.20 & 0.53 & 0.09 & 3.41 & 4.4 \\
CaseText + Previous Similar Cases & 0.16 & 0.03 & 0.19 & 0.52 & 0.08 & 3.67 & 4.7 \\
CaseText + Statutes + Precedents & 0.16 & 0.03 & 0.20 & 0.53 & 0.09 & 3.92 & 5.2 \\
CaseFacts Only & 0.15 & 0.02 & 0.17 & 0.52 & 0.08 & 3.74 & 4.6 \\
Facts + Statutes + Precedents & 0.15 & 0.02 & 0.19 & 0.52 & 0.07 & 3.08 & 4.1 \\
\bottomrule
\end{tabular}
}
\vspace*{-2mm}
\caption{Comparison of explanation generation across different legal context pipelines}
% The best result has been marked in bold.}
\label{tab:explanation_pipeline_results}
\end{table*}
%%%%%%%%%%%%%%%%%%%%%%%%%%%%%%%%%%%%%%%%%%%%%%%%%%%%